% Section 6.13 — Segmentation Robustness / Failure Analysis
% Robustness analysis for segmentation.

\section{Segmentation Robustness and Failure Analysis}
\label{sec:segmentation-robustness}

Segmentation robustness analysis documents known failure modes and limitations of the segmentation models. These experiments were conducted by Nguyen Van Phong, and the findings are reproduced from the accepted evidence.

\subsection{Known Failure Modes}
\label{sec:segmentation-failures}

The following failure modes have been documented in the retained segmentation evidence:

\begin{itemize}
    \item False-positive disease masks on Healthy images, where the segmentation model produces spurious BG or WSSV polygons despite the absence of disease.
    \item Under-segmentation of WSSV lesions on specimens with mild or early-stage infection.
    \item Boundary ambiguity between BG and WSSV regions in co-infection cases.
\end{itemize}

\subsection{Operational Implications}
\label{sec:segmentation-operational}

The false-positive failure mode on Healthy images motivates the classification-gated cascade architecture described in Section~\ref{sec:segmentation-leakage}. However, the cascade also introduces a failure mode of its own: a diseased image that is misclassified as Healthy by the classifier will suppress segmentation entirely, potentially missing a detectable disease. End-to-end validation of the cascade is therefore necessary before any operational deployment.